\section{固有性与赋值判据}\label{src:4.8}
\begin{epigraph}
处理非紧空间, 就像试图把零钱留在有破洞的口袋里.
\epigraphauthor{Angelo Vistoli}
\end{epigraph}
在定义代数叠态射的普遍闭性, 分离性与固有性之后, 我们将证明赋值判据 (定理~\bookref{4.8.7}), 推广概形通常的赋值判据 (定理~\bookref{A.4.5}). 这些赋值判据在模理论中至关重要. 本书将用它们证明 $\overline{\mathcal M}_g$ 固有 (定理~\bookref{6.5.23}), 以及 $\cBun^{\mathrm{ss}}_{r,d}(C)$ 普遍闭.

\subsection{定义}\label{src:4.8.1-subsection}
定义代数叠态射的分离性与固有性需要稍加小心. 回顾定义~\bookref{4.3.11}, 代数叠之间的可表态射 $\mathcal X\to\mathcal Y$ 称为分离的, 若其对角线 $\mathcal X\to\mathcal X\times_{\mathcal Y}\mathcal X$ (它可由概形表出) 固有.
\begin{definition}\label{src:4.8.1}
\begin{enumerate}[(1)]
\item 代数叠的态射 $\mathcal X\to\mathcal Y$ 称为\newterm{普遍闭的}{universally closed}, 若对代数叠的每个态射 $\mathcal Y'\to\mathcal Y$, 态射 $\mathcal X\times_{\mathcal Y}\mathcal Y'\to\mathcal Y'$ 都诱导闭映射 $|\mathcal X\times_{\mathcal Y}\mathcal Y'|\to|\mathcal Y'|$.
\item 代数叠的可表态射 $\mathcal X\to\mathcal Y$ 称为\newterm{固有的}{proper}, 若它普遍闭, 分离且有限型.
\item 代数叠的态射 $\mathcal X\to\mathcal Y$ 称为\newterm{分离的}{separated}, 若可表态射 $\mathcal X\to\mathcal X\times_{\mathcal Y}\mathcal X$ 固有.
\item 代数叠的态射 $\mathcal X\to\mathcal Y$ 称为固有的, 若它普遍闭, 分离且有限型.
\end{enumerate}
\end{definition}
对于概形的态射, 固有性等价于对角线为闭浸入. 对代数空间也成立, 因为代数空间的固有单态射是闭浸入 (推论~\bookref{5.5.14}). 对代数叠的态射则不成立, 因为对角线未必是单态射. 回顾域值点 $x:\Spec k\to\mathcal X$ 的稳定子 $G_x$ 由下列 Cartesian 图给出:
\[
\begin{tikzcd}
G_x\arrow[r]\arrow[d]&\Spec k\arrow[d,"{(x,x)}"]\\
\mathcal X\arrow[r]\arrow[ur,phantom,"\square" description]&\mathcal X\times\mathcal X.
\end{tikzcd}
\]
若 $\mathcal X$ 是概形 $S$ 上的分离代数叠, 则 $G_x$ 是 $k$ 上的固有群代数空间, 而由定理~\bookref{5.5.28}, 它甚至是群概形. 若 $\mathcal X$ 分离且具有仿射对角线, 则 $G_x$ 固有且仿射, 因而有限.
\begin{example}\label{src:4.8.2}
若 $G$ 是抽象有限群, 则 $BG\to\Spec\mathbb Z$ 固有. 若 $G\to S$ 是固有群概形 (例如 Abel 概形), 则 $B\mathbb G_m\to S$ 固有. 另一方面, 若 $G\to S$ 是非有限的仿射群概形 (例如 $\mathbb G_m\to\Spec\mathbb Z$), 则 $BG\to S$ 普遍闭但不分离.
\end{example}
\begin{exercise}[简单]\label{src:4.8.3}
证明 $\mathcal M_{(0,0)}$, $\mathcal M_{(0,1)}$, $\mathcal M_{(0,2)}$ 与 $\mathcal M_{(1,0)}$ 均不分离.
\end{exercise}
\begin{example}[$\cBun_{r,d}(C)$ 不分离]\label{src:4.8.4}
由于 $\cBun_{r,d}(C)$ 具有仿射对角线 (例~\bookref{4.3.15}), 且自同构群无限, $\cBun_{r,d}(C)$ 在基域 $k$ 上不分离.
\end{example}
\begin{remark}\label{src:4.8.5}
由于普遍闭性是靶上的光滑局部性质, $\mathcal X\to\mathcal Y$ 普遍闭当且仅当对从概形出发的所有映射 $T\to\mathcal Y$, $\mathcal X\times_{\mathcal Y}T\to T$ 都是闭映射; 等价地, 对一个光滑呈示 $V\to\mathcal Y$, $\mathcal X\times_{\mathcal Y}V\to V$ 普遍闭. 在定理~\bookref{4.8.7} 的证明中, 我们将说明, 当 $\mathcal Y$ 为 Noether 叠时, 只需考虑从概形出发的有限型基变换 $T\to\mathcal Y$.
\end{remark}
\newpage
\begin{exercise}\label{src:4.8.6}
域 $k$ 上的代数群 $G$ 在代数空间 $U$ 上的作用称为\newterm{固有作用}{proper action}, 若作用映射
\[
\Psi:G\times U\longrightarrow U\times U,\qquad(g,u)\longmapsto(gu,u)
\]
固有.
\begin{enumerate}[(a)]
\item 证明 $G$ 在 $U$ 上的作用固有当且仅当 $[U/G]$ 分离.
\item 对 $u\in U(k)$, 设 $\Psi_u:G\to U$ 是由 $g\mapsto gu$ 定义的映射 (通过第二个分量上的投影, 将 $\Psi$ 视为 $U$ 上的态射, 则 $\Psi_u$ 就是 $\Psi$ 在 $u$ 上的纤维). 证明以下条件等价:
\begin{enumerate}[(i)]
\item $\Psi_u:G\to U$ 固有;
\item $u:\Spec k\to[U/G]$ 固有;
\item $Gu\subseteq U$ 闭且 $G_u$ 固有.
\end{enumerate}
\end{enumerate}
\textit{(b) 的提示:} 为证明 (i) 或 (ii) 蕴含 (iii), 将 $U$ 替换为约化轨道 $(Gu)_{\mathrm{red}}$, 用一般平坦性~\bookeqref{4.3.32} 证明态射 $\Phi_u$ 或 $u$ 忠实平坦, 然后用 fppf 下降. 另见 \cite[引理 $0.3$]{GIT}.
\end{exercise}

\subsection{赋值判据}\label{src:4.8.2-subsection}
对于模问题, 固有性的赋值判据可表述为一个几何问题: 离散赋值环 $R$ 的分式域上的对象 (我们将它视为一条去掉一点的曲线), 是否可以唯一地延拓为该离散赋值环 (或曲线) 上的族. 在下面的表述中, 我们将使用注记~\bookref{4.7.3} 中正式定义的 2-交换图的提升及其态射.
\begin{theorem}[固有性/普遍闭性/分离性的赋值判据]\label{src:4.8.7}
设 $f:\mathcal X\to\mathcal Y$ 是局部 Noether 代数叠之间的拟紧且拟分离态射. 考虑 2-交换图
\begin{equation}\label{src:4.8.8}
\begin{tikzcd}[column sep=3em,row sep=3em]
\Spec K\arrow[r]\arrow[d]&\mathcal X\arrow[d,"f"]\\
\Spec R\arrow[r]&\mathcal Y.
\arrow[Rightarrow,from=1-2,to=2-1,shorten <=1.7em,shorten >=1.7em,"\alpha"']
\end{tikzcd}
\end{equation}
其中 $R$ 是分式域为 $K$ 的离散赋值环. 则
\begin{enumerate}[(1)]
\item $f$ 普遍闭当且仅当对于每个图~\bookeqref{4.8.8}, 存在离散赋值环的扩张 $R\to R'$, 使 $K\to K'=\operatorname{Frac}(R)$ 的超越次数有限, 并存在提升
\begin{equation}\label{src:4.8.9}
\begin{tikzcd}[column sep=2.4em]
\Spec K'\arrow[r]\arrow[d]&\Spec K\arrow[r]\arrow[d]&\mathcal X\arrow[d,"f"]\\
\Spec R'\arrow[r]\arrow[urr,dashed]&\Spec R\arrow[r]&\mathcal Y;
\end{tikzcd}
\end{equation}
\item $f$ 固有当且仅当 $f$ 有限型, 且对每个图~\bookeqref{4.8.8}, 存在离散赋值环的扩张 $R\to R'$, 使分式域上的映射 $K\to K'$ 的超越次数有限, 并存在如图~\bookeqref{4.8.9} 的提升, 而且该提升在相差唯一同构的意义下唯一;
\item $f$ 分离当且仅当图~\bookeqref{4.8.8} 的任意两个提升之间存在唯一同构;
\item $f$ 具有分离对角线当且仅当图~\bookeqref{4.8.8} 的一个提升的每个自同构均为平凡的.
\end{enumerate}
\end{theorem}
在下一小节证明此定理之前, 我们先给出一些例子.
\newpage
\begin{exercise}[值得练习]\label{src:4.8.10}
用赋值判据证明以下结论:
\begin{enumerate}[(a)]
\item 若 $G$ 是抽象有限群, 则 $BG\to\Spec\mathbb Z$ 固有.
\item 映射 $B\mathbb G_m\to\Spec\mathbb Z$ 普遍闭但不分离.
\end{enumerate}
\end{exercise}
\begin{exercise}[困难]\label{src:4.8.11}
证明习题~\bookref{4.1.21}(b) 中引入的稳定椭圆曲线叠 $\overline{\mathcal M}_{1,1}$ 在 $\Spec\mathbb Z$ 上固有.

我们将在后面证明, 只要 $g\geq2$, 或 $g=1$ 且 $n\geq1$, 或 $g=0$ 且 $n\geq3$, $\overline{\mathcal M}_{g,n}$ 就固有 (定理~\bookref{6.5.23}).
\end{exercise}
\begin{exercise}[值得练习]\label{src:4.8.12}
证明当 $g\geq2$ 时, $\mathcal M_g$ 在 $\Spec\mathbb Z$ 上分离.
\end{exercise}
\begin{example}[$\overline{\mathcal M}_{g,n}$ 的固有性]\label{src:4.8.13}
只要 $2g-2+n>0$ (即排除 $(g,n)=(0,0)$, $(0,1)$, $(0,2)$ 与 $(1,0)$), 模叠 $\overline{\mathcal M}_{g,n}$ 就在 $\Spec\mathbb Z$ 上固有 (定理~\bookref{6.5.23}). 我们将在 \S\bookref{6.5.1-subsection} 中, 通过验证赋值判据, 在特征 $0$ 的情形下证明这一点.
\end{example}
\begin{exercise}[值得练习]\label{src:4.8.14}
对于域 $k$ 上的光滑, 几何连通射影曲线 $C$, 证明代数叠 $\Coh_{r,d}(C)$ 与 $\cBun_{r,d}(C)$ 满足普遍闭性的赋值判据. 注意在这种情形下, 赋值判据无需扩张离散赋值环即可成立.
\end{exercise}
\begin{remark}\label{src:4.8.15}
注意上面的习题并不蕴含 $\cBun_{r,d}(C)$ 普遍闭. 事实上, 普遍闭性的赋值判据要求拟紧性, 而 $\cBun_{r,d}(C)$ 不拟紧 (见习题~\bookref{9.2.3}). 实际上, 概形之间的普遍闭态射必然拟紧 \cite[Tag 04XU]{SP}, 而对代数叠的态射想必也成立. 曲线上的向量丛叠与所有曲线组成的叠都不是普遍闭的 (见习题~\bookref{4.8.16} 与~\bookref{6.4.21}).
\end{remark}
\begin{exercise}\label{src:4.8.16}
\begin{enumerate}[(a)]
\item 证明 $\cBun_{r,d}(C)$ 不是普遍闭的.
\item 对于 $k$ 上任意维数的整固有概形 $X$, 证明 $X$ 上的无挠层叠 $\Coh^{\mathrm{tf}}$ 满足普遍闭性的赋值判据.
\end{enumerate}
\end{exercise}
\textbf{是否必须作基变换?} 对概形态射 $X\to Y$ 的固有性赋值判据~\bookeqref{A.4.5}, 无需允许扩张离散赋值环, 即存在唯一提升
\[
\begin{tikzcd}
\Spec K\arrow[r]\arrow[d]&X\arrow[d]\\
\Spec R\arrow[r]\arrow[ur,dashed]&Y.
\end{tikzcd}
\]
对代数空间的态射也成立 \cite[Tag 0A40]{SP}. 对一般的代数叠态射, 则必须允许扩张.
\begin{example}\label{src:4.8.17}
考虑域 $k$ 上的 Deligne--Mumford 叠 $\mathcal X=B\mu_n$, 其中 $k$ 的特征与 $n$ 互素. 设 $R=k[x]_{(x)}$, 其分式域为 $K=k(x)$. 则由 $x\mapsto x^n$ 给出的 $\Spec K\to\Spec K$ 是非平凡的主 $\mu_n$-torsor. 分类映射 $\Spec K\to B\mu_n$ 不能延拓为映射 $\Spec R\to B\mu_n$ (因为 $\Spec R$ 上的所有 $\mu_n$-torsor 都是平凡的). 但注意, 主 $\mu_n$-丛 $\Spec K\to\Spec K$ 在沿 $K$ 的域扩张 $K(x^{1/n})$ 作基变换后变为平凡, 因而复合 $\Spec K(x^{1/n})\to B\mu_n$ 延拓为分类平凡丛的映射 $\Spec R[x^{1/n}]\to B\mu_n$.
\end{example}
\begin{example}\label{src:4.8.18}
设 $X=\Spec R$ 是一个离散赋值环的谱, $\mathcal X$ 是 $n$ 次根叠 $X(\sqrt[n]{\mathcal O_X/\pi})$ (例~\bookref{4.9.22}). 态射 $\mathcal X\to X$ 在一般点上是同构, 但截面 $\Spec K\to\mathcal X$ 不能延拓为全局截面 $X\to\mathcal X$.
\end{example}
若 $G$ 是域上的\newterm{特殊代数群}{special algebraic group} (即每个主 $G$-丛都 Zariski 局部平凡), 例如 $\operatorname{SL}_n$ 或 $\operatorname{GL}_n$, 则 $BG$ 无需基变换就满足普遍闭性的赋值判据: 每个映射 $\Spec K\to BG$ 对应于平凡主 $G$-丛, 因而可延拓为映射 $\Spec R\to BG$. 另一方面, 对 $B\operatorname{PGL}_n$ 则必须作基变换.
\begin{exercise}[困难]\label{src:4.8.19}
证明存在一个离散赋值环的分式域上的主 $\operatorname{PGL}_n$-丛, 它不能延拓到该离散赋值环上.
\end{exercise}
对于 $\overline{\mathcal M}_g$ 的固有性赋值判据, 必须扩张离散赋值环 (见例~\bookref{6.5.11}). 另一方面, $\cBun_{r,d}(C)$ 的普遍闭性赋值判据无需扩张即可成立 (习题~\bookref{4.8.14}).

\subsection{赋值判据的证明}\label{src:4.8.3-subsection}
我们修改 \S\bookref{A.4} 中给出的概形赋值判据的证明. 出发点是闭态射的提升判据, 它推广引理~\bookref{A.4.1}.
\begin{lemma}\label{src:4.8.20}
设 $f:\mathcal X\to\mathcal Y$ 是代数叠之间的拟紧态射. 则 $f$ 是闭映射当且仅当对每个点 $x\in|\mathcal X|$, 每个特殊化 $f(x)\rightsquigarrow y_0$ 都可以提升为特殊化 $x\rightsquigarrow x_0$.
\end{lemma}
\begin{proof}
若 $f$ 是闭映射, 则 $f(\overline{\{x\}})$ 闭且等于 $\overline{\{f(x)\}}$. 因而每个点 $y_0\in\overline{\{f(x)\}}=f(\overline{\{x\}})$ 都是 $x$ 的一个特殊化的像. 对于反方向, 将 $\mathcal X$ 替换为任意闭子叠后, 只需证明 $f(|\mathcal X|)$ 闭. 由习题~\bookref{4.3.31}(a), 只需证明 $f(|\mathcal X|)$ 对特殊化封闭.
\end{proof}
与概形的情形一样, 特殊化由从离散赋值环出发的映射诱导 (命题~\bookref{A.4.4}).
\begin{proposition}\label{src:4.8.21}
若 $f:\mathcal X\to\mathcal Y$ 是 Noether 代数叠之间的有限型态射, $x\in|\mathcal X|$, 且 $f(x)\rightsquigarrow y_0$ 是特殊化, 则存在图
\[
\begin{tikzcd}
\Spec K\arrow[r]\arrow[d]&\mathcal X\arrow[d,"f"]\\
\Spec R\arrow[r]&\mathcal Y
\end{tikzcd}
\qquad
\begin{tikzcd}[column sep=2.4em]
x\arrow[d,mapsto]&{}\\
f(x)\arrow[r,rightsquigarrow]&y_0,
\end{tikzcd}
\]
其中 $R$ 是分式域为 $K$ 的离散赋值环, $\Spec K\to\mathcal X$ 的像为 $x$, 而 $\Spec R\to\mathcal Y$ 实现特殊化 $f(x)\rightsquigarrow y_0$. 特别地, Noether 代数叠中的每个特殊化 $x\rightsquigarrow x_0$, 都可由从离散赋值环出发的映射 $\Spec R\to\mathcal X$ 实现.
\end{proposition}
\begin{proof}
设 $V\to\mathcal Y$ 是光滑呈示, $v_0\in V$ 是 $y_0$ 的一个原像. 由于 $V\to\mathcal Y$ 光滑, 它是开态射 (习题~\bookref{4.3.27}), 因而在 $f(x)\rightsquigarrow y_0$ 上存在特殊化 $v\rightsquigarrow v_0$ (习题~\bookref{4.3.31}(b)). 设 $x'\in|\mathcal X_V|$ 是 $v\in V$ 与 $x\in|\mathcal X|$ 的共同原像. 设 $U\to\mathcal X_V$ 是光滑呈示, $u\in U$ 是 $x'$ 的一个原像. 将命题~\bookref{A.4.4} 应用于概形态射 $U\to V$ 及 $u\mapsto v\rightsquigarrow v_0$, 就得到所需的图.
\end{proof}
\begin{proof}[定理~\bookref{4.8.7} 的证明]
$f:\mathcal X\to\mathcal Y$ 的拟分离性意味着对角线 $\Delta_f:\mathcal X\to\mathcal X\times_{\mathcal Y}\mathcal X$ 与双重对角线 $\Delta_{\Delta_f}:\mathcal X\to\mathcal X\times_{\mathcal X\times_{\mathcal Y}\mathcal X}\mathcal X$ 都拟紧. 在证明中, 由赋值判据推出普遍闭性, 分离性与对角线分离性时, 分别需要 $f$, $\Delta_f$ 与 $\Delta_{\Delta_f}$ 的拟紧性.

我们先建立普遍闭性的判据. 假设赋值判据成立, 但 $f:\mathcal X\to\mathcal Y$ 不是普遍闭的. 由于普遍闭性是靶上的光滑局部性质, 可假设 $\mathcal Y=Y$ 是 Noether 概形, 且存在概形之间的映射 $T\to Y$, 使 $f_T:\mathcal X_T\to T$ 不是闭映射. 我们将归约到 $T\to Y$ 为有限型态射的情形. 由引理~\bookref{4.8.20}, 存在 $z\in|\mathcal X_T|$ 及特殊化 $f_T(z)\rightsquigarrow t_0$, 而它不能提升为特殊化 $z\rightsquigarrow z_0$. 这意味着 $\mathcal Z=\overline{\{z\}}\subseteq\mathcal X_T$ 与纤维 $(\mathcal X_T)_{t_0}$ 不相交. 若 $p:X\to\mathcal X$ 是光滑呈示, 则 $\mathcal Z$ 在 $X_T\to\mathcal X_T$ 下的原像 $Z$ 不与纤维 $(X_T)_{t_0}$ 相交. 将引理~\bookref{A.4.8} 应用于 $X_T\to T$, 可知将 $T$ 替换为 $t_0$ 的一个开邻域之后, 存在分解 $T\xrightarrow{g}T'\to Y$ 及闭子概形 $Z'\subseteq X_{T'}$, 使 $T'\to Y$ 有限型, $Z\cap(X_{T'})_{g(t_0)}=\varnothing$, 且 $\operatorname{im}(Z\hookrightarrow X_T\to X_{T'})\subseteq Z'$. 设 $z'\in|\mathcal X_{T'}|$ 是 $z\in|\mathcal X_T|$ 的像, 则 $z'$ 映到 $g(f_T(z))\in T'$, 而特殊化 $g(f_T(z))\rightsquigarrow g(t_0)$ 不能提升为 $z'$ 的特殊化. 由引理~\bookref{4.8.20}, $\mathcal X_{T'}\to T'$ 不是闭映射.

若 $T\to Y$ 是有限型态射, 则基变换 $\mathcal X_T\to T$ 是 Noether 代数叠之间的有限型态射, 并且也满足普遍闭性的赋值判据. 因而只需证明 $f:\mathcal X\to\mathcal Y$ 是闭映射. 由引理~\bookref{4.8.20}, 需要证明, 给定点 $x\in|\mathcal X|$, 每个特殊化 $f(x)\rightsquigarrow y_0$ 都可提升为特殊化 $x\rightsquigarrow x_0$. 由命题~\bookref{4.8.21}, 存在图
\begin{equation}\label{src:4.8.22}
\begin{tikzcd}
\Spec K\arrow[r]\arrow[d]&\mathcal X\arrow[d,"f"]\\
\Spec R\arrow[r]&\mathcal Y
\end{tikzcd}
\qquad
\begin{tikzcd}[column sep=2.4em]
x\arrow[d,mapsto]&{}\\
f(x)\arrow[r,rightsquigarrow]&y_0.
\end{tikzcd}
\end{equation}
\begin{samepage}
赋值判据表明, 将 $R$ 替换为一个扩张之后, 存在提升 $\Spec R\to\mathcal X$, 它进而给出提升 $f(x)\rightsquigarrow y_0$ 的特殊化 $x\rightsquigarrow x_0$.
\par
\end{samepage}

反过来, 假设 $f:\mathcal X\to\mathcal Y$ 普遍闭, 且给定如图~\bookeqref{4.8.22} 的图. 将 $\mathcal Y$ 替换为 $\Spec R$, 将 $\mathcal X$ 替换为 $\mathcal X\times_{\mathcal Y}\Spec R$ 后, 可假设 $\mathcal Y=\Spec R$, 且有图
\[
\begin{tikzcd}
\Spec K\arrow[r]\arrow[dr]&\mathcal X\arrow[d,"f"]\\
{}&\Spec R.
\end{tikzcd}
\]
将 $\mathcal X$ 替换为 $\overline{\{x\}}$ 后, 可假设 $\mathcal X$ 整且一般点为 $x$. 由于 $\mathcal X\to\Spec R$ 是闭映射, 存在特殊化 $x\rightsquigarrow x_0$, 它映到 $\Spec R$ 中从一般点到闭点的特殊化. 由于 $f$ 拟分离, $\Spec K\to\mathcal X$ 拟紧. 将命题~\bookref{4.8.21} 应用于 $\Spec K\to\mathcal X$, 得到分式域为 $K'$ 的离散赋值环 $R'$ 及交换图
\[
\begin{tikzcd}[column sep=2.5em]
\Spec K'\arrow[r]\arrow[d]&\Spec K\arrow[d]\arrow[dr]&{}\\
\Spec R'\arrow[r]&\mathcal X\arrow[r,"f"]&\Spec R,
\end{tikzcd}
\]
使 $\Spec R'\to\mathcal X$ 实现特殊化 $x\rightsquigarrow x_0$. 由于 $\Spec R'\to\Spec R$ 满射, $R\to R'$ 是离散赋值环的扩张, 而 $\Spec R'\to\mathcal X$ 给出图~\bookeqref{4.8.22} 的一个提升.

对角线分离的判据留给读者 (习题~\bookref{4.8.23}). 假设 $f$ 具有分离对角线, 则 $f:\mathcal X\to\mathcal Y$ 的分离性赋值判据转化为对角线 $\Delta_f:\mathcal X\to\mathcal X\times_{\mathcal Y}\mathcal X$ 的普遍闭性赋值判据. 因而固有性与分离性的赋值判据由普遍闭性的赋值判据推出. 另见 \cite[定理 $7.10$]{LMB00}, \cite[\S$4$]{Fal03}, 以及 \cite[Tags 0CLQ, 0CLS, 0CLV, 0CLY]{SP}.
\end{proof}
\begin{exercise}\label{src:4.8.23}
验证对角线分离的赋值判据.

试用类似于对角线非分歧的无穷小提升判据 (习题~\bookref{4.7.5}) 的论证, 改用 $\Delta_f$ 分离当且仅当 $\Delta_{\Delta_f}$ 是闭浸入这一事实 (习题~\bookref{4.3.38}), 并注意, 由于 $\Delta_{\Delta_f}$ 是可由概形表出的拟紧单态射, 可以用概形的赋值判据~\bookeqref{A.4.5} 来验证它是闭浸入.
\end{exercise}
